% ---------------------------------------------------------------------------
% APPLICATION / IMPLEMENTATION SECTION — contribution to the joint report.
% To merge into the master document: delete the preamble and the
% \begin{document} ... \end{document} wrapper, paste the \section block below,
% and fold the thebibliography entries into the master bibliography.
% ---------------------------------------------------------------------------
\documentclass[11pt,a4paper]{article}

\usepackage[utf8]{inputenc}
\usepackage[T1]{fontenc}
\usepackage{lmodern}
\usepackage[margin=2.1cm]{geometry}
\usepackage{booktabs}
\usepackage{tabularx}
\usepackage{listings}
\usepackage{xcolor}
\usepackage{float}
\usepackage{hyperref}
\usepackage{microtype}

\hypersetup{colorlinks=true, linkcolor=black, urlcolor=blue, citecolor=black}

\lstset{
  basicstyle=\ttfamily\scriptsize,
  numbers=left,
  numberstyle=\tiny\color{gray},
  breaklines=true,
  frame=single,
  framesep=3pt,
  keywordstyle=\color{blue!70!black},
  stringstyle=\color{red!60!black},
  showstringspaces=false,
  captionpos=b,
  aboveskip=6pt,
  belowskip=4pt
}

\setlength{\parskip}{0.2em}

\begin{document}

\section{Application Implementation}

The preceding sections established that running a full Large Language Model on consumer
mobile hardware is impractical, and that a client--server architecture --- the device acting
as a user-facing interface to a remotely hosted model --- is the most consistent way to
deliver an LLM-powered experience on mobile. This section documents the realisation of that
architecture in a working application.

\emph{MotivBot} administers the well-being battery described above and returns individualised
supportive feedback generated by a remote language model. The participant first describes
their emotional state in free text, then completes the 62-item questionnaire, after which both
the narrative and the structured responses are transmitted to the model. Turkish and Persian
are supported, the latter requiring right-to-left layout. The application was built twice:
natively for iOS in Swift, and subsequently for Android in Kotlin.

\subsection{The Two Native Languages}

Apple relied on Objective-C until Swift was introduced at the 2014 Worldwide Developers
Conference. Swift was designed to eliminate the characteristics of Objective-C that most often
produced defects, and three are directly relevant here \cite{swiftbook}. \textbf{Optionality}
requires that any variable which may hold no value be declared with the \texttt{?} suffix, and
the compiler refuses to read it unwrapped, eliminating a class of null-dereference crashes
before the program runs. \textbf{Value semantics} mean a \texttt{struct} is copied on
assignment rather than shared, so passing a response to another component cannot let that
component mutate the original. \textbf{Structured concurrency} via \texttt{async}/\texttt{await}
allows a network call to be written sequentially, replacing the nested completion handlers that
previously fragmented control flow. The interface uses SwiftUI, Apple's declarative framework,
in which the developer describes what the interface should contain for a given state rather
than issuing instructions to mutate it. Development requires Xcode on macOS; the application
targets iOS 26.2.

Android development was historically conducted in Java. Kotlin, developed by JetBrains, became
an official Android language in 2017 and Google's preferred language in 2019 \cite{kotlin}. Its
design goals closely parallel Swift's, which is what made the port tractable: nullability is
expressed in the type system in the same way, the \texttt{data class} serves the role of the
Swift \texttt{struct}, and coroutines provide concurrency, with \texttt{suspend} functions
corresponding to Swift's \texttt{async} functions. The interface uses Jetpack Compose
\cite{compose}, close enough to SwiftUI in philosophy that most screens translated structurally
rather than being redesigned. The application targets API 35, supports API 26 and above, and
compiles against JVM 17.

\begin{table}[H]
\centering
\small
\begin{tabularx}{\textwidth}{@{}lXX@{}}
\toprule
\textbf{Concern} & \textbf{iOS (Swift)} & \textbf{Android (Kotlin)} \\
\midrule
Null safety & Optionals (\texttt{?}, \texttt{guard let}) & Nullable types (\texttt{?}, \texttt{?.}) \\
Data modelling & \texttt{struct}, \texttt{Codable} & \texttt{data class}, \texttt{@Serializable} \\
Concurrency & \texttt{async}/\texttt{await}, \texttt{Task} & \texttt{suspend}, coroutines, \texttt{Flow} \\
UI framework & SwiftUI (declarative) & Jetpack Compose (declarative) \\
State observation & \texttt{@Observable} & \texttt{StateFlow} \\
Thread guarantee & \texttt{@MainActor} & \texttt{Dispatchers.Main} \\
Toolchain & Xcode (macOS only) & Android Studio, Gradle \\
\bottomrule
\end{tabularx}
\caption{Correspondence between the two native language ecosystems.}
\end{table}

\subsection{Cross-Platform Alternatives}

The survey above examined Flutter, React Native and Kotlin Multiplatform, observing that
cross-platform frameworks reduce cost at the price of indirect hardware access. That analysis
preceded implementation; what building the application revealed refines it.

The iOS implementation comprises 13 Swift files totalling 2{,}059 lines; the Android
implementation, 35 Kotlin files totalling 4{,}020 lines, an increase that overstates duplicated
effort since much of it reflects explicit dependency injection and a separated data layer
rather than restated logic. The more serious cost was not the writing but the
\textbf{divergence}: several defects were corrected in the Kotlin implementation --- most
consequentially, that the model was being instructed in English regardless of interface
language --- and none propagated to the Swift implementation automatically. Two codebases
holding the same intent separate under maintenance unless deliberate effort keeps them aligned,
and that effort is a permanent tax rather than a one-time cost.

Kotlin Multiplatform occupies a position between the two extremes surveyed: rather than sharing
the entire application, it shares only the layers containing no platform-specific concerns ---
here the questionnaire definitions, prompt construction and network client --- while the
interface remains fully native \cite{kmp}. This suits the application, because the layer that
most needs to remain consistent is precisely the one containing no user interface: the
instrument definitions and their scoring must be identical across platforms or the research
data is not comparable. Native development was nonetheless retained, because the iOS
application existed before the Android requirement arose, because right-to-left Persian layout
is more reliably obtained from the platform than from a framework's reimplementation of it, and
because building the same application twice yields a direct comparison of the two ecosystems.

\subsection{LLM Integration}

No model executes on the device. The application issues an HTTPS request to a remotely hosted
model and renders the response, so the memory-hierarchy constraints examined earlier --- the
windowing and bundling techniques required to run inference under a constrained DRAM budget ---
do not arise. The request uses \texttt{gpt-4o-mini} with temperature 0.8, a 500-token limit and
a 60-second timeout. The temperature is deliberately high: a low value produces near-identical
phrasing across sessions, undermining the impression of individualised feedback, while the
supportive register is constrained by the system prompt instead. Tokens are not streamed, since
the response is short and presented as a single passage.

Each request carries a system message establishing the assistant's role and a user message
containing the participant's narrative followed by every item paired with its response and
score. The most instructive defect in this project concerned language: the original
implementation sent a fixed English system prompt while questionnaire content reached the model
in the participant's language, leaving the output language undefined in practice, so a
participant reading a Turkish interface could receive an English reply. The correction was to
construct the system prompt as a function of the selected language.

\begin{lstlisting}[language=Java,caption={Language-conditioned system prompt construction.}]
private fun systemPrompt(language: AppLanguage): String {
    val script = if (language.isRtl)
        " Write in the Persian script, right-to-left, using Persian (not Arabic) digits." else ""
    return """
        You MUST write your entire response in ${language.promptName}.$script
        You are a compassionate psychological well-being assistant. ...
        The questionnaire below may reach you in ${language.promptName} or in English;
        either way, your answer must be entirely in ${language.promptName}.
    """.trimIndent()
}
\end{lstlisting}

The language instruction appears both before and after the assistant persona, because a single
instruction at the head of a prompt is followed unreliably once several thousand tokens of
questionnaire data intervene. The Persian branch also specifies the digit set, because a model
asked to write Persian will otherwise frequently emit Arabic-Indic numerals, which Persian
readers find incongruous.

\subsection{Technology Stack and Justification}

\begin{table}[H]
\centering
\small
\begin{tabularx}{\textwidth}{@{}llX@{}}
\toprule
\textbf{Concern} & \textbf{Library} & \textbf{Reason for selection} \\
\midrule
\multicolumn{3}{@{}l}{\textit{iOS}} \\
UI & SwiftUI & Declarative; first-party; no third-party dependency \\
Networking & URLSession & Built in; native \texttt{async}/\texttt{await} support \\
Serialisation & Codable & Compiler-generated; no reflection or manual parsing \\
\midrule
\multicolumn{3}{@{}l}{\textit{Android}} \\
UI & Jetpack Compose \cite{compose} & Declarative; closest structural analogue to SwiftUI \\
Networking & Retrofit + OkHttp \cite{retrofit} & Declarative HTTP interfaces; native \texttt{suspend} support \\
Serialisation & kotlinx.serialization & Compile-time generated, matching \texttt{Codable} \\
DI & Hilt \cite{hilt} & Compile-time graph validation; lifecycle-aware scoping \\
Database & Room \cite{room} & Compile-time SQL verification against the schema \\
Preferences & DataStore & Asynchronous and transactional, unlike SharedPreferences \\
Testing & JUnit, MockWebServer & Network behaviour testable without a live endpoint \\
\bottomrule
\end{tabularx}
\caption{Libraries selected for each platform, with justification.}
\end{table}

A consistent principle governs these selections: \textbf{prefer libraries that move errors from
runtime to compile time}. Retrofit validates endpoint declarations during compilation; Room
parses SQL and verifies it against the declared schema; Hilt validates the dependency graph
before the application runs; kotlinx.serialization generates serialisers at compile time rather
than resolving fields reflectively. In each case a failure that would otherwise appear on a
participant's device appears instead on the developer's screen --- and for an application
administering clinical instruments, where a runtime failure mid-assessment loses the session,
that is the correct trade-off. Two translations required more than substituting one library for
its counterpart: Swift's \texttt{URLRequest.timeoutInterval} bounds a request as a whole,
whereas OkHttp exposes connect, read and write phases separately, so equivalent behaviour was
reconstructed by giving each phase a budget and setting \texttt{callTimeout} to cap the total;
and the Swift service signalled failure by throwing a \texttt{LocalizedError}, whereas the
Kotlin data source returns a sealed \texttt{AppResult} type, so the compiler enforces that every
failure mode is handled at the call site. Correctness is verified by 33 unit tests across four
suites, covering questionnaire progression, prompt and CSV construction, the network layer
against a mock server, and input validation.

In summary, building natively twice confirmed both the expected benefit and the expected cost:
platform behaviour such as right-to-left Persian layout came without effort, while the two
codebases began to diverge as soon as defects were corrected in one and not the other, which is
why Kotlin Multiplatform is recommended for future work. The most transferable finding, however,
concerns LLM integration rather than platform choice: the language defect described above passed
every conventional test --- the code compiled, the request succeeded, and a well-formed response
was returned; only its language was wrong. Applications built on language models therefore
require a form of verification that conventional testing does not supply, directed at whether
the model was asked the right question rather than at whether the request was correctly sent.

\begin{thebibliography}{9}
\footnotesize
\setlength{\itemsep}{0pt}
\setlength{\parskip}{0pt}

\bibitem{swiftbook} Apple Inc., \emph{The Swift Programming Language}. [Online]. Available:
\url{https://docs.swift.org/swift-book/}

\bibitem{kotlin} JetBrains, \emph{Kotlin Documentation}. [Online]. Available:
\url{https://kotlinlang.org/docs/home.html}

\bibitem{compose} Google, \emph{Jetpack Compose}. [Online]. Available:
\url{https://developer.android.com/compose}

\bibitem{kmp} JetBrains, \emph{Kotlin Multiplatform}. [Online]. Available:
\url{https://kotlinlang.org/docs/multiplatform.html}

\bibitem{retrofit} J. Wharton et al., ``Retrofit: A type-safe HTTP client for Android and
Java,'' Square. [Online]. \url{https://square.github.io/retrofit/}

\bibitem{hilt} Google, \emph{Dependency Injection with Hilt}. [Online]. Available:
\url{https://developer.android.com/training/dependency-injection/hilt-android}

\bibitem{room} Google, \emph{Room Persistence Library}. [Online]. Available:
\url{https://developer.android.com/training/data-storage/room}

\end{thebibliography}

\end{document}
